\documentclass{article}
\usepackage{bm}
\usepackage{mathtools}
\usepackage{amsthm}
\usepackage{fancyhdr}
\usepackage{float}
\usepackage{listings}
\lstset{breaklines=true}
\usepackage{amssymb}
\usepackage{braket}
\usepackage{graphicx}
\usepackage{xcolor}
\usepackage{titlesec}
\usepackage{titling}
\usepackage{tabularx}
\usepackage{enumitem}
\usepackage{amsmath}
\usepackage{amsfonts}
\usepackage{hyperref}
\usepackage[margin=1in]{geometry}

\hypersetup{
	colorlinks=true,
	linkcolor=blue,
	filecolor=magenta,
	urlcolor=blue,
	citecolor=blue,
	anchorcolor=blue
}

\renewcommand{\thesubsubsection}{\Roman{subsubsection}}
\title{Math 3IA3 - Assignment 4}
\author{Matthew Farah, \href{mailto:farahm11@mcmaster.ca}{$\langle$farahm11@mcmaster.ca$\rangle$}, 400468251}
\date{\today}

\makeatletter
\setlist{listparindent=\parindent}

\pagestyle{fancy}
\fancyhead{}
\fancyhead[L]{Matthew Farah, $\langle$\href{mailto:farahm11@mcmaster.ca}{farahm11@mcmaster.ca}$\rangle$, 400468251}
\fancyhead[R]{Math 3IA3 - Assignment 4}

\newcolumntype{Y}{>{\centering\arraybackslash}X}
\setlength{\parindent}{0cm}

\setcounter{section}{4}

\newcounter{chapter}
\renewcommand{\thechapter}{\arabic{chapter}}
\setcounter{chapter}{2}

\newenvironment{chaptertitle}[1][]{
	\newpage
	\refstepcounter{chapter}
	\begin{center}
	\vspace*{.5em}
	\par
	\Large
	Section \thechapter:\hspace{-.5em}
	\addcontentsline{toc}{section}{\protect{\large{Section \thechapter $\,$ (#1)}}}
}{
	\end{center}
	\vspace{2.5em}
}

\newcounter{exercise}
\renewcommand{\theexercise}{\arabic{exercise}}

\newenvironment{exercise}[1][]{
	\refstepcounter{exercise}
	\addcontentsline{toc}{subsection}{\protect{\large{Exercise \thesection.\thechapter.#1}}}
	\vspace{1em}
	\par
	\large\textbf{Exercise \thesection.\thechapter.#1}
	\vspace{.2cm}
	\par
	\setcounter{subexercise}{0}
	\setcounter{step}{0}
}{\vspace{.5em}}

\newenvironment{solution}{
	\vspace{1em}\par\large\textbf{{Solution:}}\par\vspace{1em}
}{
	\vspace{1em}
	\newpage
}

\newcounter{subexercise}
\renewcommand{\thesubexercise}{\alph{subexercise}}

\newenvironment{subexercise}{
	\refstepcounter{subexercise}
	\vspace{1em}\par\large\textbf{(\thesubexercise)}\hspace{-.5em}
	\setcounter{step}{0}
	\addcontentsline{toc}{subsubsection}{\protect{\large{Part \theexercise.\thesubexercise}}}
}{
	\par
	\vspace{1em}
}

\makeatother

\makeatletter

\newcounter{step}
\renewcommand{\thestep}{\Roman{step}}

\newenvironment{step}{
	\refstepcounter{step}
	\vspace{1.5em}\par\noindent\large\textbf{(\thestep)}
}{
	\par
	\vspace{1.5em}
}

\makeatother

\newcolumntype{Y}{>{\centering\arraybackslash}X}
\renewcommand{\arraystretch}{1.8}

\fontsize{10pt}{14pt}\selectfont

\begin{document}

\maketitle
\pagestyle{fancy}
\cfoot{\thepage}

\tableofcontents
\newpage

\begin{exercise}[12]
	Find functions $f$ and $g$ defined on $\left( 0, \infty \right)$ such that $\lim_{x\to\infty}f = \infty$ and $\lim_{x\to\infty}g = \infty$ and $\lim_{x\to\infty}\left( f - g \right) = 0$. Can you find such functions, with $g(x) > 0$ for all $x > 0$ for all $x \in \left( 0, \infty \right)$, such that $\lim_{x\to\infty} \frac{f}{g} = 0$?
\end{exercise}

\begin{solution}

\begin{step}
Find $f$ and $g$ such that $\lim_{x\to\infty}(f-g)=0$ while both diverge.
\end{step}

A suitable choice is
\[
f(x)=x+\frac{1}{x},\qquad g(x)=x.
\]

\begin{align*}
\lim_{x\to\infty}f(x) &= \lim_{x\to\infty}\left(x+\frac{1}{x}\right) = \infty,\\
\lim_{x\to\infty}g(x) &= \lim_{x\to\infty}x = \infty,\\
\lim_{x\to\infty}(f(x)-g(x)) &= \lim_{x\to\infty}\frac{1}{x} = 0.
\end{align*}

\begin{step}
Find $f$ and $g$ with $g(x)>0$ and $\displaystyle\lim_{x\to\infty}\frac{f}{g}=0$.
\end{step}

Let
\[
f(x)=\sqrt{x},\qquad g(x)=x.
\]

\begin{align*}
\lim_{x\to\infty}f(x) &= \infty,\\
\lim_{x\to\infty}g(x) &= \infty,\qquad g(x)>0\ \forall x>0,\\
\lim_{x\to\infty}\frac{f(x)}{g(x)} &= \lim_{x\to\infty}\frac{\sqrt{x}}{x}
= \lim_{x\to\infty}\frac{1}{\sqrt{x}} = 0.
\end{align*}

\end{solution}

\setcounter{section}{5}
\setcounter{chapter}{1}

\begin{exercise}[12]
	Suppose that $f\colon \mathbb{R} \rightarrow \mathbb{R}$ is continuous on $\mathbb{R}$ and that $f(r) = 0$ for every rational number $r$. Prove that $f(x) = 0$ for all $x \in \mathbb{R}$.
\end{exercise}

\begin{solution}

Assume, for contradiction, that there exists $k\in\mathbb{R}$ with $f(k)\neq 0$.  
Let $L=f(k)$, so $|L|>0$.

Since $f$ is continuous at $k$, for  
\[
\varepsilon=\frac{|L|}{2}>0,
\]
there exists $\delta>0$ such that
\[
0<|x-k|<\delta \implies |f(x)-L|<\varepsilon.
\]

Because $\mathbb{Q}$ is dense in $\mathbb{R}$, there exists a rational number $r$ with  
\[
0<|r-k|<\delta.
\]
Then
\[
|f(r)-L| = |0-L| = |L| \ge \frac{|L|}{2} = \varepsilon,
\]
contradicting continuity at $k$.

Thus no such $k$ exists and $f(x)=0$ for all real $x$.

\end{solution}

\setcounter{chapter}{3}

\begin{exercise}[12]
	Let $I \coloneq \left[0, \frac{\pi}{2} \right]$ and let $f\colon I \rightarrow \mathbb{R}$ be defined by $f(x) \coloneq \sup{\set{x^2, \cos(x)}}$ for $x \in I$. Show that there exists an absolute maximum point $x_0 \in I$ for $f$ on $I$. Show that $x_0$ is a solution to the equation $\cos(x) = x^2$.
\end{exercise}

\begin{solution}

\begin{step}
Show that $f$ attains an absolute maximum on $I$.
\end{step}

Both $x^2$ and $\cos x$ are continuous on $I$. Thus
\[
f(x)=\max\{x^2,\cos x\}
\]
is continuous as the maximum of two continuous functions. Since $I$ is compact,
$f$ attains an absolute maximum at some $x_0\in I$.

\begin{step}
Show that any interior maximizer satisfies $x_0^2=\cos x_0$.
\end{step}

Suppose $x_0\in(0,\frac{\pi}{2})$ is a point where $f$ achieves its maximum.

\textbf{Case 1:} $\cos(x_0)>x_0^2$.
	Then in a small neighbourhood of $x_0$ we have $f(x)=\cos x$. Since $\cos x$ is strictly decreasing on $[0,\frac{\pi}{2}]$, choosing $x<x_0$ sufficiently close gives $\cos(x)>\cos(x_0)$, contradicting maximality.

\textbf{Case 2:} $x_0^2>\cos(x_0)$.
Then near $x_0$ we have $f(x)=x^2$. Since $x^2$ is strictly increasing on $[0,\frac{\pi}{2}]$, choosing $x>x_0$ sufficiently close gives $x^2>x_0^2$, again contradicting maximality.

Thus neither case is possible, so
\[
x_0^2 = \cos(x_0).
\]

\end{solution}

\setcounter{chapter}{6}

\begin{exercise}[7]
	Let $I \subseteq \mathbb{R}$ be an interval and let $f\colon I \rightarrow \mathbb{R}$ be increasing on $I$. If $c$ is not an endpoint of $I$, show that the jump $j_{f}(c)$ of $f$ at $c$ is given as
	\[
	\inf\{f(y)-f(x) : x<c<y,\; x,y\in I\}.
	\]
\end{exercise}

\begin{solution}

Let
\[
j_f(c)=f(c+)-f(c-),
\]
where $f(c+)=\lim_{y\downarrow c}f(y)$ and $f(c-)=\lim_{x\uparrow c}f(x)$ exist because $f$ is increasing.

Let
\[
S=\{\,f(y)-f(x):x<c<y,\ x,y\in I\,\}.
\]

\begin{step}
Show $j_f(c)$ is a lower bound of $S$.
\end{step}

If $x<c<y$, then by monotonicity,
\[
f(y)-f(x)\ge f(c+)-f(c-)=j_f(c).
\]

\begin{step}
Show $\inf S \le j_f(c)$.
\end{step}

Let $\varepsilon>0$.  
Choose $y>c$ close enough so that
\[
f(y)<f(c+)+\varepsilon/2,
\]
and choose $x<c$ close enough so that
\[
f(x)>f(c-)-\varepsilon/2.
\]
Then
\begin{align*}
f(y)-f(x)
&< \big(f(c+)+\varepsilon/2\big) - \big(f(c-)-\varepsilon/2\big) \\
&= j_f(c)+\varepsilon.
\end{align*}

Thus $\inf S \le j_f(c)+\varepsilon$ for all $\varepsilon>0$, so $\inf S \le j_f(c)$.

Combining both steps gives $\inf S = j_f(c)$.

\end{solution}

\setcounter{section}{6}
\setcounter{chapter}{1}


\begin{exercise}[9]
	Show that if $f$ is an even function and has a derivative at every point, then the derivative $f'$ is an odd function. Also prove that if $g \colon \mathbb{R} \rightarrow \mathbb{R}$ is a differentiable and odd function then $g'$ is an even function.
\end{exercise}

\begin{solution}

\begin{step}
If $f$ is even and differentiable, then $f'$ is odd.
\end{step}

Since $f$ is even, $f(-x)=f(x)$ for all $x$.

For any $x$,
\begin{align*}
f'(-x)
&= \lim_{h\to 0}\frac{f(-x+h)-f(-x)}{h} \\
&= \lim_{h\to 0}\frac{f(x-h)-f(x)}{h} \\
&= \lim_{k\to 0}\frac{f(x+k)-f(x)}{-k} \qquad (k=-h)\\
&= -f'(x).
\end{align*}

Thus $f'$ is odd.

\begin{step}
If $g$ is odd and differentiable, then $g'$ is even.
\end{step}

Since $g$ is odd, $g(-x)=-g(x)$.

For any $x$,
\begin{align*}
g'(-x)
&= \lim_{h\to 0}\frac{g(-x+h)-g(-x)}{h} \\
&= \lim_{h\to 0}\frac{-g(x-h)+g(x)}{h} \\
&= \lim_{k\to 0}\frac{g(x+k)-g(x)}{k} \qquad (k=-h)\\
&= g'(x).
\end{align*}

Thus $g'$ is even.

\end{solution}

\end{document}
